% !TeX root = ../master.tex

\lecture{4}{2026-09-15}{Conduction, Pt. III}

\begin{problem}{4--17}
  Metal plates ($k = \qty{180}{\W\per\m\per\K}, \rho = \qty{2800}{\kg\per\m\cubed}$, and $c_p = \qty{880}{\J\per\kg\per\K}$) with a thickness of \qty{1}{\cm} are being heated in an oven for 2 minutes.
  Air in the oven is maintained at \qty{800}{\celsius} with a convection heat transfer coefficient of \qty{200}{\W\per\m\squared\per\K}.
  If the initial temperature of the plates is \qty{20}{\celsius}, determine the temperature of the plates when they are removed from the oven.
\end{problem}

\begin{problemwithimage}{p4_1}{4--34}
  Carbon steel balls ($\rho = \qty{7833}{\kg\per\m\cubed}, k = \qty{54}{\W\per\m\per\K}, c_p = \qty{0,465}{\kJ\per\kg\per\celsius}$, and $\alpha = \qty{1,474e-6}{\m\squared\per\s}$) \qty{8}{\mm} in diameter are annealed by heating them first to \qty{900}{\celsius} in a furnace and then allowing them to cool slowly to \qty{100}{\celsius} in ambient air at \qty{35}{\celsius}.
  If the average heat transfer coefficient is \qty{75}{\W\per\m\squared\per\K}, determine how long the annealing process will take.
  If \num{2500} balls are to be annealed per hour, determine the total rate of heat transfer from the balls to the ambient air.
\end{problemwithimage}

\begin{problemwithimage}{p4_2}{4--47}
  A hot brass plate is having its upper surface cooled by impinging jet of air at temperature of \qty{15}{\celsius} and convection heat transfer coefficient of \qty{220}{\W\per\m\squared\per\K}.
  The \qty{10}{\cm} thick brass plate ($\rho = \qty{8530}{\kg\per\m\cubed}, c_p = \qty{380}{\J\per\kg\per\K}, k = \qty{110}{\W\per\m\per\K}$, and $\alpha = \qty{33,9e-6}{\m\squared\per\s}$ has a uniform initial temperature of \qty{650}{\celsius}, and the bottom surface of the plate is insulated.
  Determine the temperature at the center plane of the brass plate after 3 minutes of cooling.
  Solve this problem using analytical one-term approximation (not the Heisler charts).
\end{problemwithimage}

\begin{problemwithimage}{p4_3}{4-66}
  A father and son conducted the following simple experiment on a hot dog which measured \qty{12,5}{\cm} in length and \qty{2,2}{\cm} in diameter.
  They inserted one food thermometer into the midpoint of the hot dog and another one was placed just under the skin of the hot dog.
  The temperature of the thermometers were monitored until both temperatures read \qty{20}{\celsius}  which is the ambient temperature.
  The hot dog was then placed in \qty{94}{\celsius} boiling water and after exactly 2 minutes they recorded the center temperature and the skin temperature of the hot dog to be \qty{59}{\celsius} and \qty{88}{\celsius}, respectively.
  Assuming the following properties for the hot god: $\rho = \qty{980}{\kg\per\m\cubed}$ and $c_p = \qty{3900}{\J\per\kg\per\K} $ and using transient temperature charts, determine
  \begin{subproblems}
    \subproblem the thermal diffusivity of the hot dog,
    \subproblem the thermal conductivity of the hot dog, and
    \subproblem the convection heat transfer coefficient.
  \end{subproblems}
\end{problemwithimage}
